% TOP_PROBLEM: {"id": 3014, "title": "Kirby Problem 5.7", "category_id": 7, "category": {"id": 7, "name": "topology", "display_name": "Topology", "description": "Properties preserved under continuous deformations.", "slug": "topology", "order_index": 7, "created_at": "2026-07-31T15:26:25.671Z"}, "status": "open", "problem_number": "KP-5.7", "set_id": 11}
% FIRST_POSTED: 2026-10-01
% ENTRY_KIND: statement_only
% UnsolvedMath ID 3014; source catalogue code KP-5.7
% !TeX program = lualatex
% Definition and sources only; this file is not an LLM solution attempt.
\documentclass[11pt,a4paper]{article}
\usepackage[margin=25mm]{geometry}
\usepackage{fontspec}
\setmainfont{lmroman10-regular.otf}[BoldFont=lmroman10-bold.otf,
 ItalicFont=lmroman10-italic.otf,BoldItalicFont=lmroman10-bolditalic.otf]
\usepackage{amsmath,amssymb,mathtools}
\usepackage{unicode-math}
\setmathfont{latinmodern-math.otf}
\AtBeginDocument{\let\setminus\smallsetminus}
\usepackage{xurl,hyperref}
\hypersetup{colorlinks=true,urlcolor=blue}
\setcounter{secnumdepth}{0}
\setlength{\parindent}{0pt}
\setlength{\parskip}{5pt}
\setlength{\emergencystretch}{3em}
\begin{document}
\section*{Kirby Problem 5.7}\label{3014}
{\small UnsolvedMath ID 3014 · KP-5.7 · Topology · Kirby's Problems in Low-Dimensional Topology}
\subsection{Definitions and mathematical statement}
Let $S^5\subset\mathbb R^6$ be the standard smooth sphere.
A smooth circle action is a smooth map
$a:S^1\times S^5\to S^5$ satisfying the identity and
multiplication laws. The stabilizer of a point $x$ is
$G_x=\{u\in S^1:a(u,x)=x\}$. An orbit is multiple or
exceptional if $G_x$ is a nontrivial finite subgroup of the
circle, hence a cyclic group of some order $m>1$.

The action is pseudo-free if there are no fixed points of the
whole circle and the exceptional orbits are isolated: each has
an invariant neighborhood containing no other exceptional orbit.
Use effective actions, so no nonidentity element fixes all points.
Away from the exceptional orbits the action is free, with
trivial stabilizers.

Does there exist such a smooth pseudo-free action on $S^5$
with more than three exceptional orbits? Equivalently, is four
or more possible in this dimension? Different exceptional orbits
are counted separately, even if their stabilizers have the same
order. The source poses the existence question underlying the
Montgomery--Yang problem.

The sphere must have its standard smooth structure. An action
on a homology sphere or another $5$-manifold would not suffice
without identifying the manifold with $S^5$. Actions with fixed
points or positive-dimensional families of exceptional orbits
are outside the pseudo-free hypothesis; their orbit counts do
not answer the question.
\subsection{Short English statement}
Can a smooth pseudo-free circle action on the standard five-sphere have more than three exceptional orbits?
\subsection{Sources}
\begin{itemize}
\item[S1] ULAM.AI and the UnsolvedMath Contributors, supplied problems.json, record 3014 (KP-5.7), Kirby's Problems in Low-Dimensional Topology. Catalogue identity and source transcription; CC BY 4.0.
\url{https://huggingface.co/datasets/ulamai/UnsolvedMath}.
\item[S2] K3: A New Problem List in Low-Dimensional Topology, Problem 5.7. Expanded formulation of the supplied catalogue statement.
\url{https://math.berkeley.edu/sites/default/files/surv-295-ruberman-watermarked-author-pdf.pdf}.
\end{itemize}
\end{document}
